\documentclass[11pt,a4paper]{article}

\usepackage[margin=2.5cm]{geometry}
\usepackage{graphicx}
\usepackage{amsmath,amssymb,amsfonts}
\usepackage{hyperref}
\usepackage{booktabs}
\usepackage{xcolor}
\usepackage{natbib}
\usepackage{authblk}
\usepackage{placeins}
\urlstyle{same}

\title{\textbf{Parameterization-Dependent Decoupling Windows in Statistical Manifolds:\\
An Empirical Study via Finite Ising Models}}

\author{Wei Ye\thanks{ORCID: \href{https://orcid.org/0009-0004-9737-3969}{0009-0004-9737-3969}}}
\affil{Independent Researcher}

\date{\today}

\begin{document}

\maketitle

\begin{abstract}
\noindent
We test whether ``decoupling windows''---parameter intervals in which the Fisher information matrix (FIM) condition number decreases monotonically---are intrinsic geometric features of a statistical manifold or artifacts of the chosen coordinates.  Using exact enumeration of $N=16$ Ising models on six graph topologies, we compare the natural coordinates $P_1=(\beta,h)$ with the canonical exponential-family coordinates $P_2=(\beta,B=\beta h)$ along a weak-field path $h=0.01$.  In $P_1$, all five non-complete topologies exhibit a positive-to-negative zero crossing of $\Delta=-d\log\kappa(\mathcal F)/d\beta$, with mean boundary $\beta_G^*=0.4789\pm0.0258$.  In $P_2$, only three of the five non-complete topologies exhibit such a crossing, and the observed boundaries move to $\beta_G^*=0.2431\pm0.0113$.  A zero-field exact-enumeration robustness check gives the same crossing pattern and nearly identical values (mean $1.1\times10^{-3}$ and maximum $1.9\times10^{-3}$ among the valid $P_2$ crossings), showing that the weak-field path is not driving the main effect.  By contrast, the scalar curvature $R(\beta)$, computed independently in both coordinate systems on the periodic Grid $4\times4$, agrees to median relative error $4.1\times10^{-4}$ and remains smooth over the scanned range.  The results provide a controlled counterexample to interpreting $\kappa(\mathcal F)$-based windows as Riemannian invariants.  They also show that previously observed graph-topology correlations of $\beta_G^*$ are not stable enough, at this system size, to support an invariant topological law.
\end{abstract}

\section{Introduction}
\label{sec:intro}

The geometry of statistical manifolds, initiated by Rao and developed in modern information geometry by Amari and Nagaoka \citep{rao1945information, amari2000methods}, equips a parametric family $p(x;\theta)$ with the Fisher information matrix (FIM)
\begin{equation}
\mathcal F_{ij}(\theta)=\mathbb E[\partial_i\log p\,\partial_j\log p],
\end{equation}
which is a Riemannian metric under regularity assumptions.  Intrinsic quantities such as scalar curvature and geodesic distance are invariant under smooth coordinate changes.  In statistical mechanics, information-geometric curvature has been studied as a marker of criticality and finite-size rounding \citep{rup1995riemannian, brody2003information, janke2004information}.

The FIM condition number
\begin{equation}
\kappa(\mathcal F)=\lambda_{\max}(\mathcal F)/\lambda_{\min}(\mathcal F)
\end{equation}
is different.  It is an eigenvalue ratio of a coordinate representation of the metric, not a scalar invariant.  Under a reparameterization $\theta\mapsto\theta'$, the metric transforms tensorially as $\mathcal F'=J^\top \mathcal F J$, but the Euclidean eigenvalues of the representing matrix are preserved only under orthogonal transformations.  This distinction is familiar in optimization, where sharpness-aware methods and sharpness-like quantities are useful operational diagnostics while remaining sensitive to the chosen coordinates \citep{foret2021sharpness, dinh2017sharp}.  FIM spectra in high-dimensional models show similar coordinate and scaling pathologies \citep{karakida2021pathological}.

This paper examines a proposed ``decoupling window'' diagnostic defined from the FIM condition number.  Let
\begin{equation}
I_{\rm eff}(\beta)=\log\kappa(\mathcal F(\beta)),\qquad
\Delta(\beta)=-\frac{d I_{\rm eff}}{d\beta}.
\end{equation}
A window is the region where $\Delta>0$, and its right boundary is operationally defined as the first positive-to-negative crossing $\beta_G^*$ of $\Delta$.  If this window were an intrinsic feature of the statistical manifold, its existence and location should not depend strongly on the coordinate system used to write the same probability family.

We test this directly in finite Ising models.  Exact enumeration at $N=16$ removes Monte Carlo thermalization and estimator noise, making it suitable for a coordinate-dependence control experiment.  The result is negative but useful: the window defined by $\kappa(\mathcal F)$ is parameterization-conditioned.  The same weak-field path gives a window in all five non-complete topologies in $P_1=(\beta,h)$, but only three of five in $P_2=(\beta,B=\beta h)$, and the observed crossing locations differ by roughly a factor of two.  An invariant scalar-curvature calculation verifies that our coordinate transformations are implemented correctly and that the contrast is not a numerical artifact.

\section{Theory}
\label{sec:theory}

\subsection{Condition Number Is Not a Scalar}

For a smooth reparameterization $\theta=\theta(\theta')$, let $J^i_{\;a}=\partial\theta^i/\partial\theta'^a$.  The FIM transforms as
\begin{equation}
\mathcal F'_{ab}=J^i_{\;a}\mathcal F_{ij}J^j_{\;b}.
\end{equation}
This is the correct tensor transformation law.  However, the matrix eigenvalues of $\mathcal F'$ depend on the coordinate basis.  Therefore $\kappa(\mathcal F)$, $I_{\rm eff}$, and $\Delta$ are not Riemannian scalars.

Within a fixed parameterization, $\Delta$ can be written using the Hellmann--Feynman derivative of the extreme eigenvalues:
\begin{equation}
\Delta
=\operatorname{tr}(\Omega\,\dot{\mathcal F}),\qquad
\Omega=\frac{v_{\min}v_{\min}^{\top}}{\lambda_{\min}}
-\frac{v_{\max}v_{\max}^{\top}}{\lambda_{\max}},
\label{eq:delta_operator}
\end{equation}
where the dot denotes differentiation along the chosen path.  Equation~\eqref{eq:delta_operator} is useful diagnostically in a fixed coordinate system, but $\Omega$, $\dot{\mathcal F}$, and their contraction are not invariant under general reparameterizations.

\subsection{Invariant Benchmark}

The scalar curvature $R$ of the Fisher metric is a coordinate invariant.  For the two-dimensional Ising model in the thermodynamic limit, information-geometric arguments predict critical scaling of curvature near the critical point \citep{janke2004information}.  Finite systems round such singularities \citep{brody2003information}; therefore the relevant use of $R$ at $N=16$ is not to estimate the thermodynamic singularity, but to verify coordinate invariance and to check whether any curvature feature is visible at the $\kappa$-window locations.

\section{Methods}
\label{sec:methods}

\subsection{Finite Ising Models}

We study Ising spins $s_i\in\{\pm1\}$ on graphs with $N=16$ vertices and Hamiltonian
\begin{equation}
\mathcal H(s)=-J\sum_{\langle i,j\rangle}s_i s_j-h\sum_i s_i,
\qquad J=1.
\end{equation}
The Boltzmann weight is
\begin{equation}
p(s\mid\beta,h)=\frac{\exp[-\beta\mathcal H(s)]}{Z(\beta,h)}
=\frac{\exp[\beta(E(s)+hM(s))]}{Z(\beta,h)},
\end{equation}
where $E(s)=\sum_{\langle i,j\rangle}s_i s_j$ and $M(s)=\sum_i s_i$.  We enumerate all $2^{16}=65{,}536$ states exactly.

The six graph topologies are: periodic Grid $4\times4$, Ring, Complete, Star, Random Regular of degree 4, and Watts--Strogatz Small World with $k=4$ and rewiring probability $p=0.3$ \citep{watts1998collective}.  The Complete graph is reported in per-topology tables but excluded from the non-complete summary statistics because its spectral gap and edge count make it an extreme outlier relative to the other five graphs.

\subsection{Two Parameterizations}

We compare two coordinate systems on the same two-dimensional statistical manifold.

\paragraph{$P_1=(\beta,h)$.}
The score statistics are
\begin{equation}
\partial_\beta\log p=(E+hM)-\langle E+hM\rangle,\qquad
\partial_h\log p=\beta(M-\langle M\rangle).
\end{equation}
Thus
\begin{equation}
\mathcal F^{(P_1)}
=\begin{pmatrix}
\operatorname{Var}(E+hM) & \beta\,\operatorname{Cov}(E+hM,M)\\
\beta\,\operatorname{Cov}(E+hM,M) & \beta^2\operatorname{Var}(M)
\end{pmatrix}.
\end{equation}

\paragraph{$P_2=(\beta,B)$ with $B=\beta h$.}
In canonical coordinates,
\begin{equation}
p(s\mid\beta,B)=\frac{\exp[\beta E(s)+B M(s)]}{Z(\beta,B)},
\end{equation}
so
\begin{equation}
\mathcal F^{(P_2)}
=\begin{pmatrix}
\operatorname{Var}(E) & \operatorname{Cov}(E,M)\\
\operatorname{Cov}(E,M) & \operatorname{Var}(M)
\end{pmatrix}.
\end{equation}
Along the weak-field path used for the main scan, $h=0.01$ and $B=\beta h$.  Along the zero-field robustness path, $h=0$ and $B=0$.  The $P_2$ metric is not degenerate at $B=0$: $\operatorname{Var}(M)$ remains finite in the finite system, although symmetry makes $\operatorname{Cov}(E,M)=0$ at exactly zero field.

\subsection{Window Extraction}

For each topology and parameterization, we scan $\beta\in[0.05,0.55]$ on 401 evenly spaced points.  We compute $\log\kappa(\mathcal F)$ and estimate $\Delta=-d\log\kappa/d\beta$ by centered differences.  The boundary $\beta_G^*$ is the first positive-to-negative crossing of $\Delta$, linearly interpolated between adjacent grid points.  If no such crossing occurs in the scanned range, $\beta_G^*$ is reported as absent.

\subsection{Scalar Curvature}

For the periodic Grid $4\times4$, we compute the scalar curvature of the two-dimensional Fisher metric using a 9-point centered stencil.  In $P_1$, the stencil steps are $\delta\beta=2\times10^{-3}$ and $\delta h=10^{-4}$.  In $P_2$, they are $\delta\beta=2\times10^{-3}$ and $\delta B=5\times10^{-5}$.  Christoffel symbols and the Riemann tensor are assembled from the first and second derivatives of the metric by their standard formulas.  Agreement between the two coordinate computations is used as an implementation check.

\section{Results}
\label{sec:results}

\subsection{Weak-Field Main Scan}
\label{subsec:window}

Table~\ref{tab:per_topology} shows the exact crossing locations at $h=0.01$.

\begin{table}[ht]
\centering
\caption{Exact-enumeration decoupling-window boundary $\beta_G^*$ at $h=0.01$.  A dash means no positive-to-negative crossing of $\Delta$ was found in $\beta\in[0.05,0.55]$.}
\label{tab:per_topology}
\begin{tabular}{lccc}
\toprule
Topology & $\lambda_2(L)$ & $P_1=(\beta,h)$ & $P_2=(\beta,B)$\\
\midrule
Grid $4\times4$ & 2.0000 & 0.4766 & 0.2309\\
Ring & 0.1522 & 0.5213 & --\\
Complete & 16.0000 & 0.2328 & 0.1450\\
Star & 1.0000 & 0.4398 & --\\
Random Regular (degree 4) & 1.3089 & 0.4802 & 0.2403\\
Small World & 0.8330 & 0.4769 & 0.2580\\
\bottomrule
\end{tabular}
\end{table}

Across the five non-complete topologies, $P_1$ gives crossings in all cases with mean
\begin{equation}
\beta_G^*(P_1)=0.4789\pm0.0258.
\end{equation}
In $P_2$, only three of the five non-complete topologies cross, with mean over valid crossings
\begin{equation}
\beta_G^*(P_2)=0.2431\pm0.0113.
\end{equation}
The main result is therefore not a stable topological law for $\beta_G^*$, but the stronger and more basic fact that the existence and location of the $\kappa$-window depend on the coordinate representation of the same finite statistical manifold.

The graph spectral-gap correlation is correspondingly not robust.  In $P_1$, the Pearson correlation between $\beta_G^*$ and $\lambda_2$ over the five non-complete graphs is $r=-0.481$ with a wide Fisher-transform confidence interval $(-0.957,0.697)$.  In $P_2$, only three non-complete graphs have crossings, making any correlation estimate underpowered and not interpretable as evidence for a topological scaling law.  We therefore do not claim a universal relation between $\beta_G^*$ and $\lambda_2$.

\subsection{Zero-Field Robustness}
\label{subsec:h0}

The weak-field path $h=0.01$ is used to avoid numerical complications in finite-difference curvature stencils and to provide a smooth two-coordinate diagnostic path.  To test whether it biases the condition-number result, we repeated the exact enumeration at $h=0$.  The crossing pattern is unchanged: $P_1$ crosses in all five non-complete topologies and $P_2$ crosses in the same three.  Table~\ref{tab:h0_check} summarizes the non-complete means.

\begin{table}[ht]
\centering
\caption{Zero-field robustness check.  Means and standard deviations are over non-complete topologies with valid crossings.}
\label{tab:h0_check}
\begin{tabular}{lcccc}
\toprule
Path & Parameterization & Mean $\beta_G^*$ & Std & Crossings\\
\midrule
$h=0.01$ & $P_1=(\beta,h)$ & 0.4789 & 0.0258 & 5/5\\
$h=0$ & $P_1=(\beta,h)$ & 0.4788 & 0.0259 & 5/5\\
$h=0.01$ & $P_2=(\beta,B)$ & 0.2431 & 0.0113 & 3/5\\
$h=0$ & $P_2=(\beta,B)$ & 0.2441 & 0.0106 & 3/5\\
\bottomrule
\end{tabular}
\end{table}

Thus the weak-field path is not a substitute for the zero-field Ising critical point, and we do not interpret it as such.  It is a numerical probe of the finite statistical manifold.  The $h=0$ check shows that the observed parameterization dependence of $\kappa(\mathcal F)$ is not an artifact of choosing $h=0.01$.

\subsection{Scalar Curvature Benchmark}
\label{subsec:curvature}

Figure~\ref{fig:curvature} shows the scalar curvature on the periodic Grid $4\times4$ along $h=0.01$.  The two independent computations, in $P_1$ and $P_2$, agree to median relative error $4.1\times10^{-4}$, with 95th percentile $1.3\times10^{-2}$.  This verifies that the stencil implementation is computing an invariant quantity.

The curvature is smooth at this system size.  Around the Onsager infinite-volume reference $\beta_c=\frac12\log(1+\sqrt2)\approx0.4407$ \citep{onsager1944crystal}, the neighboring sampled values give $R\approx0.237$ (with linear interpolation changing this only at the displayed precision), and no divergence is visible.  It also shows no feature at the $P_1$ window near $\beta\approx0.477$ or the $P_2$ window near $\beta\approx0.23$.  This is consistent with strong finite-size rounding at $N=16$ and should not be read as a thermodynamic-limit statement about $R$.

\section{MCMC Finite-Size Scaling Attempt}
\label{sec:mcmc}

Exact enumeration is limited to small $N$.  A natural next step would be Monte Carlo finite-size scaling of the invariant scalar curvature.  We attempted this with a 9-point stencil and MCMC-estimated FIM elements.  The pilot results are not used as positive evidence for a curvature signal; they are included to document why the direct approach is statistically unattractive.

If a metric element has relative Monte Carlo error $\sigma_g/g\sim N_{\rm eff}^{-1/2}$, a centered second derivative with stencil width $\delta$ amplifies the dominant error by approximately $\delta^{-2}$.  Thus
\begin{equation}
\frac{\sigma_R}{|R|}
\sim \frac{1}{\sqrt{N_{\rm eff}}\delta^2},
\qquad
N_{\rm eff}\gtrsim \delta^{-4}
\end{equation}
is required for order-one relative error, up to model-dependent prefactors.  For $\delta=2\times10^{-3}$, this scaling gives $N_{\rm eff}$ of order $6.25\times10^{10}$ per stencil point before accounting for autocorrelation near criticality.  The recorded $L=4$ pilot values fluctuate in sign and magnitude, as shown in Fig.~\ref{fig:R_noise}, matching the expectation that the finite-difference curvature estimate is noise dominated.

This is not a proof that all MCMC approaches to $R$ are impossible.  Reweighting, analytic cumulant formulas, cluster algorithms, or variance-reduction methods may improve the situation.  The claim here is narrower: naive finite-difference recovery of scalar curvature from independently estimated MCMC FIM elements is not a reliable line of evidence for this manuscript.

\section{Discussion}
\label{sec:discussion}

The central lesson is that a condition-number window should be labeled with its coordinate system.  A more precise notation is $\beta_{G,P}^*$ or ``the $P$-window'', where $P$ denotes the parameterization.  Without that qualifier, $\kappa(\mathcal F)$ can be mistaken for an intrinsic geometric observable.

The result also changes how topology should be discussed.  The corrected degree-4 Random Regular graph and the exact zero-field check do not support the earlier strong claim that the $\beta_G^*$--$\lambda_2$ relation has an invariant sign or universal slope.  What remains is a methodological negative result: topology correlations inferred from $\kappa(\mathcal F)$ are conditional on the coordinate representation and on the operational definition of the crossing.

This does not make $\kappa(\mathcal F)$ useless.  As a coordinate diagnostic, the FIM condition number can still summarize anisotropic parameter sensitivity, identify ill-conditioned directions, and support comparisons made within a fixed parameterization.  This is the same pragmatic status held by many sharpness-like diagnostics in optimization: they can guide practice without being intrinsic geometric observables \citep{dinh2017sharp, foret2021sharpness}.  What fails here is the stronger invariant interpretation.  A $\kappa(\mathcal F)$ window should not be promoted to a Riemannian feature of the statistical manifold unless it is replaced by or connected to scalar quantities.

\section{Conclusion}
\label{sec:conclusion}

Exact enumeration of finite Ising models shows that decoupling windows defined by decreasing FIM condition number are parameterization-conditioned.  Along the same weak-field path, $P_1=(\beta,h)$ gives crossings in all five non-complete topologies near $\beta\approx0.479$, whereas $P_2=(\beta,B)$ gives crossings in only three topologies near $\beta\approx0.243$.  A zero-field exact check reproduces the same pattern, so the effect is not caused by the weak field.  The scalar curvature benchmark behaves as an invariant and remains smooth at $N=16$.

The main contribution is therefore a controlled falsification of an invariance claim: $\kappa(\mathcal F)$-based decoupling windows may be useful coordinate diagnostics, but the finite systems studied here show that they are not intrinsic geometric features of the statistical manifold.  Establishing whether a useful invariant counterpart exists remains open and will require either analytic curvature control or lower-noise estimators than the naive finite-difference MCMC route tested here.

\section*{Code and data availability}

Code, data, and reproduction scripts are available in the public GitHub repository release \href{https://github.com/OZYW/fisher-ising-decoupling-windows/releases/tag/v1.0}{OZYW/fisher-ising-decoupling-windows v1.0}.

\section*{Use of AI-assisted tools}

The author used AI-assisted tools during exploratory hypothesis generation, code drafting and debugging, numerical cross-checking, manuscript revision, and review of reproducibility materials.  The research direction, selection of tests, interpretation of results, verification of numerical outputs, reference checking, and final manuscript decisions were performed by the author, who takes full responsibility for the content, claims, code, and citations.

\begin{figure}[ht]
\centering
\includegraphics[width=0.95\textwidth]{figure1_exact_decoupling_windows.png}
\caption{Exact decoupling-window scan at $h=0.01$ with zero-field robustness summary.  Top row: $I_{\rm eff}=\log\kappa(\mathcal F)$ in $P_1=(\beta,h)$ and $P_2=(\beta,B)$.  Middle row: $\Delta=-dI_{\rm eff}/d\beta$ with positive-to-negative crossings marked.  Bottom left: $\beta_G^*$ versus Laplacian spectral gap for the weak-field path.  Bottom right: comparison of weak-field and zero-field exact crossings; circles are $P_1$ and squares are $P_2$.}
\label{fig:exact_windows}
\end{figure}

\begin{figure}[ht]
\centering
\includegraphics[width=0.88\textwidth]{figure2_scalar_curvature_grid4x4.png}
\caption{Scalar curvature $R(\beta)$ on the periodic Grid $4\times4$ at $h=0.01$, computed independently in $P_1=(\beta,h)$ and $P_2=(\beta,B)$.  The two coordinate computations agree to median relative error $4.1\times10^{-4}$.  $R$ is smooth at $N=16$ and shows no feature at the parameterization-dependent condition-number windows.}
\label{fig:curvature}
\end{figure}

\begin{figure}[ht]
\centering
\includegraphics[width=0.88\textwidth]{figure3_mcmc_noise_diagnostic.png}
\caption{Recorded $L=4$ pilot MCMC scalar-curvature estimates and finite-difference SNR scaling.  The sign-flipping values are used only as a diagnostic of noise amplification in naive stencil-based MCMC curvature recovery.}
\label{fig:R_noise}
\end{figure}

\FloatBarrier

\bibliographystyle{plainnat}
\begin{thebibliography}{99}

\bibitem[Amari and Nagaoka(2000)]{amari2000methods}
S.~Amari and H.~Nagaoka.
\newblock \textit{Methods of Information Geometry}.
\newblock American Mathematical Society, 2000.

\bibitem[Brody and Ritz(2003)]{brody2003information}
D.~C. Brody and A.~Ritz.
\newblock Information geometry of finite Ising models.
\newblock \textit{Journal of Geometry and Physics}, 47(2--3):207--220, 2003.

\bibitem[Dinh et~al.(2017)]{dinh2017sharp}
L.~Dinh, R.~Pascanu, S.~Bengio, and Y.~Bengio.
\newblock Sharp minima can generalize for deep nets.
\newblock In \textit{Proceedings of the 34th International Conference on Machine Learning}, PMLR 70:1019--1028, 2017.

\bibitem[Foret et~al.(2021)]{foret2021sharpness}
P.~Foret, A.~Kleiner, H.~Mobahi, and B.~Neyshabur.
\newblock Sharpness-aware minimization for efficiently improving generalization.
\newblock In \textit{International Conference on Learning Representations}, 2021.

\bibitem[Janke et~al.(2004)]{janke2004information}
W.~Janke, D.~A. Johnston, and R.~Kenna.
\newblock Information geometry and phase transitions.
\newblock \textit{Physica A: Statistical Mechanics and its Applications}, 336(1--2):181--186, 2004.

\bibitem[Karakida et~al.(2021)]{karakida2021pathological}
R.~Karakida, S.~Akaho, and S.~Amari.
\newblock Pathological spectra of the Fisher information metric and its variants in deep neural networks.
\newblock \textit{Neural Computation}, 33(8):2274--2307, 2021.

\bibitem[Onsager(1944)]{onsager1944crystal}
L.~Onsager.
\newblock Crystal statistics. I. A two-dimensional model with an order-disorder transition.
\newblock \textit{Physical Review}, 65(3--4):117--149, 1944.

\bibitem[Rao(1945)]{rao1945information}
C.~R. Rao.
\newblock Information and the accuracy attainable in the estimation of statistical parameters.
\newblock \textit{Bulletin of the Calcutta Mathematical Society}, 37:81--91, 1945.

\bibitem[Ruppeiner(1995)]{rup1995riemannian}
G.~Ruppeiner.
\newblock Riemannian geometry in thermodynamic fluctuation theory.
\newblock \textit{Reviews of Modern Physics}, 67:605--659, 1995.

\bibitem[Watts and Strogatz(1998)]{watts1998collective}
D.~J. Watts and S.~H. Strogatz.
\newblock Collective dynamics of ``small-world'' networks.
\newblock \textit{Nature}, 393:440--442, 1998.

\end{thebibliography}

\end{document}
